\documentclass{article}
% ----------------------------------------------------------------------
% Define external packages, language, margins, fonts, new commands 
% and colors
% ----------------------------------------------------------------------
\usepackage[utf8]{inputenc} % Codification
\usepackage[english]{babel} % Writing idiom

\usepackage{pgfplots}
\pgfplotsset{compat=1.18}
\usepackage[T1]{fontenc}
\usepackage{imakeidx}
\usepackage{circuitikz}
\usetikzlibrary{positioning}
\usepackage[export]{adjustbox} % Align images
\usepackage{amsmath} % Extra commands for math mode
\usepackage{caption}
\usepackage{amssymb} % Mathematical symbols
\usepackage{anysize} % Personalize margins
    %\marginsize{2cm}{2cm}{2cm}{2cm} % {left}{right}{above}{below}
\usepackage{appendix} % Appendices
\usepackage{cancel} % Expression cancellation
\usepackage{placeins}
\usepackage{caption} % Captions
    \captionsetup{labelfont={bf}}
\usepackage{cite} % Citations, like [1 - 3]
\usepackage{color} % Text coloring
\usepackage{fancyhdr} % Head note and footnote
    \pagestyle{fancy}
    \fancyhf{}
    \fancyhead[L]{
        \includegraphics[width=4cm]{Imagens/NovaFct.png}
    } % Left of Head notea
    \fancyhead[R]{\footnotesize Controlo em Sistemas Ciber-físicos} % Right of Head note
    \fancyfoot[L]{\footnotesize MEEC} % Left of Footnote
    \fancyfoot[R]{\thepage  } % Center of Footnote
    %\fancyfoot[R]{\footnotesize Degree} % Right of Footnote
    \renewcommand{\footrulewidth}{0.4pt} % Footnote rule
\usepackage{float} % Utilization of [H] in figures
\usepackage{graphicx} % Figures in LaTeX
\usepackage[colorlinks = true, plainpages = true, linkcolor = istblue, urlcolor = istblue, citecolor = istblue, anchorcolor = istblue]{hyperref}
\usepackage{indentfirst} % First paragraph
\usepackage[super]{nth} % Superscripts
\usepackage{siunitx} % SI units

\usepackage{tocloft}%style the listofcontent to have dots
\renewcommand{\cftsecleader}{\cftdotfill{\cftdotsep}}

\usepackage[list=true, listformat=simple]{subcaption} % Subfigures

\usepackage{titlesec} % Font
    \titleformat{\section}{\Large\bfseries}{\thesection}{1em}{}
    \titleformat{\subsection}{\large\bfseries}{\thesubsection}{1em}{}
    \titleformat{\subsubsection}{\normalsize\bfseries}{\thesubsubsection}{1em}{}
    %\fancyfoot[C]{\thepage}

%code
\usepackage[breakable]{tcolorbox}

\graphicspath{ {Imagens/.} }

% Code highlighting
\usepackage{listings}
\lstset{ 
    language=Matlab,
    frame=single,
    numbers=left,
    numberstyle=\tiny,
    stepnumber=1,
    numbersep=5pt,
    backgroundcolor=\color{cellbackground}, % background color for the code block
    keywordstyle=\color{blue}\bfseries,
    commentstyle=\color{dkgreen},
    stringstyle=\color{red},
    breaklines=true,
    breakatwhitespace=false,
    showspaces=false,
    showstringspaces=false,
    showtabs=false,
    captionpos=b,
}

%pseudocode styling
\newcounter{nalg}[section] % defines algorithm counter for chapter-level
\DeclareCaptionLabelFormat{simple}{Algoritmo \thenalg} % defines a new caption label as Algorithm x.y
\lstnewenvironment{algorithm}[1][] % defines the algorithm listing environment
{
    \vspace{1ex}
    \refstepcounter{nalg} % increments algorithm number
    \captionsetup{labelformat=simple,labelsep=colon,position=top} % set caption position to top
    \lstset{ % this is the style
        mathescape=true,
        frame=tB,
        numbers=left, 
        numberstyle=\tiny,
        basicstyle=\scriptsize, 
        keywordstyle=\color{black}\bfseries\em,
        keywords={input, output, return, datatype, function, in, if, else, foreach, while, begin, end}, % add the keywords you want, or load a language as Rubens explains in his comment above.
        numbers=left,
        xleftmargin=.04\textwidth,
        #1 % this is to add specific settings to an usage of this environment (for instance, the caption and referable label)
    }
}
{}

% Exact colors from NB
\definecolor{incolor}{HTML}{303F9F}
\definecolor{outcolor}{HTML}{D84315}
\definecolor{cellborder}{HTML}{CFCFCF}
\definecolor{cellbackground}{HTML}{F7F7F7}
\definecolor{commentcolour}{rgb}{0.42, 0.45, 0.51}
\definecolor{string}{rgb}{0.42, 0.45, 0.51}

\usetikzlibrary{shapes,arrows}
\tikzstyle{block} = [draw, rectangle, 
    minimum height=2.5em, minimum width=3.5em]
\tikzstyle{sum} = [draw, circle, node distance=1cm]
\tikzstyle{input} = [coordinate]
\tikzstyle{output} = [coordinate]
\tikzstyle{pinstyle} = [pin edge={to-,thin,black}]


% Random text (not needed)
\usepackage{lipsum}
\usepackage{duckuments}

\usepackage{cleveref}

\usepackage{xcolor}

% New and re-newcommands
\newcommand{\sen}{\operatorname{\sen}} % Sine function definition
\newcommand{\HRule}{\rule{\linewidth}{0.5mm}} % Specific rule definition
\renewcommand{\appendixpagename}{\LARGE Appendices}

% Colors
\definecolor{istblue}{RGB}{28, 118, 196}
\definecolor{dkgreen}{rgb}{0,0.6,0}

\usepackage{amsmath,amssymb}
\usepackage{bm}

%%%%%%%%%%%%%%%%%%%%%%%%%%%%%%%%%%%%%%%%%%%%%%%%%%%%%%%%%%%%%%%%%%%%%%%%
%                                 Document                             %
%%%%%%%%%%%%%%%%%%%%%%%%%%%%%%%%%%%%%%%%%%%%%%%%%%%%%%%%%%%%%%%%%%%%%%%%
\begin{document}

% ----------------------------------------------------------------------
% Cover
% ----------------------------------------------------------------------
\begin{center}
    \begin{figure}
        \vspace{-1.0cm}
        \includegraphics[scale = 1, left]{Imagens/NovaFct.png} % Nova logo
    \end{figure}

    \mbox{}\\[2.0cm]
    \textsc{\Huge Mestrado em Engenharia Eletrotécnica e de Computadores}\\[2.5cm]
    \textsc{\LARGE Controlo em Sistemas Ciber-físicos}\\[2.0cm]
    \HRule\\[0.4cm]
    {\large \bf { 
        Controlling Satellite Dynamics\\
    }[\texttt{EN}]}\\[0.2cm]
    \HRule\\[1.5cm]
\end{center}

\begin{flushleft}
    \textbf{Authors:}
\end{flushleft}

\begin{center}
    \begin{minipage}{0.5\textwidth}
        \begin{flushleft}
            Filipe Cavalheiro (70119)\\
        \end{flushleft}
    \end{minipage}%
    \begin{minipage}{0.5\textwidth}
        \begin{flushright}
            \href{mailto:fs.cavalheiro@campus.fct.unl.pt}{\texttt{fs.cavalheiro@campus.fct.unl.pt}}
        \end{flushright}
    \end{minipage}
\end{center}

    \vspace{4cm}

    \begin{center}
        \large \bf 2025/2026 -- 1º Semestre
    \end{center}

    \thispagestyle{empty}

    \setcounter{page}{0}

    \newpage

    \newpage

    \renewcommand{\contentsname}{Table of contents}
    \tableofcontents % Generates the table of contents
    \newpage

    \renewcommand{\listfigurename}{List of Images}
    \listoffigures
    \newpage

\section*{Multi-Agent Drone System on a Sphere}

Consider a set of $N$ drones constrained to move on the surface of a sphere of radius $R > 0$.

\subsection*{State Definition}

For each drone $i \in \{1,\dots,N\}$, the state vector is defined as
\begin{equation}
    \bm{x}_i =
    \begin{bmatrix}
    \bm{p}_i \\
    \bm{v}_i \\
    \mathrm{vec}(\bm{R}_i) \\
    \bm{\omega}_i
    \end{bmatrix}
    \in \mathbb{R}^{18},
\end{equation}

where:
\begin{itemize}
    \item $\bm{p}_i \in \mathbb{R}^3$ is the position,
    \item $\bm{v}_i \in \mathbb{R}^3$ is the linear velocity,
    \item $\bm{R}_i \in SO(3)$ is the rotation matrix,
    \item $\bm{\omega}_i \in \mathbb{R}^3$ is the angular velocity.
\end{itemize}

\subsection*{Geodesic Distance Between Agents}

For agents $i$ and $j$, the geodesic distance on the sphere is
\begin{equation}
    d_{ij} = R \arccos\!\left(
    \frac{\tilde{\bm{p}}_i \cdot \tilde{\bm{p}}_j}{R^2}
    \right).    
\end{equation}


\subsection*{Reference Point Initialization}

A set of $N$ reference points $\{\bm{p}_i^{\mathrm{ref}}\}$ is generated on the sphere using a Fibonacci lattice:
\[
\|\bm{p}_i^{\mathrm{ref}}\| = R.
\]

The minimum safe distance is chosen as the average nearest-neighbor distance:
\begin{equation}
    d_{\min} = \frac{1}{N} \sum_{i=1}^N
    \frac{1}{N-1} \sum_{\substack{j=1 \\ j \neq i}}^N d_{ij}^{\mathrm{ref}}.
\end{equation}

$\frac{1}{N-1}$ because of $j \neq i$


\subsection*{Projection onto the Sphere}

At each time step $k$, the corresponding unit normal vector is: 
\[
\bm{n}_i = \frac{\bm{p}_i}{\|\bm{p}_i\|}.
\]


the position of each drone is projected onto the sphere:
\[
\tilde{\bm{p}}_i = R \bm{n}_i.
\]

\subsection*{Lyapunov-Based Distance Control}

Define the distance error
\[
e_{ij} = d_{ij} - d_{\min}.
\]

The control interaction is active only when agents are closer than the safe distance:
\[
\text{if } e_{ij} \le 0.
\]

The pairwise repulsion direction is approximated by the Euclidean direction
\[
\bm{d}_{ij} = \frac{\tilde{\bm{p}}_i - \tilde{\bm{p}}_j}
{\|\tilde{\bm{p}}_i - \tilde{\bm{p}}_j\|}.
\]

The cumulative control vector for agent $i$ is
\begin{equation}
    \bm{u}_i =
    \begin{cases}
    - k \displaystyle\sum_{\substack{j=1 \\ j \neq i}}^N e_{ij} \bm{d}_{ij}, & \text{if } e_{ij} \le 0, \\[6pt]
    - k \displaystyle\sum_{\substack{j=1 \\ j \neq i}}^N 0, & \text{if } e_{ij} > 0.
    \end{cases}
\end{equation}
where $k > 0$ is a gain.

\subsection*{Projection onto the Tangent Space}

To enforce motion on the sphere, the control is projected onto the tangent plane:
\[
\bm{u}_i^{\perp} =
\bm{u}_i - (\bm{u}_i^\top \bm{n}_i)\bm{n}_i.
\]

\subsection*{Tracking Controller}

Each drone tracks the reference position $\bm{p}_i^{d}$ using a geometric controller
(Mellinger-type), producing thrust and torque inputs:
\[
(\bm{T}_i, \bm{\tau}_i)
= \mathcal{C}(\bm{p}_i, \bm{v}_i, \bm{R}_i, \bm{\omega}_i,
\bm{p}_i^{d}).
\]

The position tracking error $\bm{e}_{p,i}$ is integrated as
\[
\bm{\xi}_{p,i}(k+1) =
\bm{\xi}_{p,i}(k) + \Delta t\, \bm{e}_{p,i}(k).
\]

\subsection*{Drone Dynamics}

The continuous-time dynamics are given by
\begin{align*}
\dot{\bm{p}}_i &= \bm{v}_i, \\
\dot{\bm{v}}_i &= \bm{f}(\bm{R}_i, \bm{T}_i), \\
\dot{\bm{R}}_i &= \bm{R}_i [\bm{\omega}_i]_\times, \\
\dot{\bm{\omega}}_i &= \bm{g}(\bm{\omega}_i, \bm{\tau}_i),
\end{align*}
where $[\cdot]_\times$ denotes the skew-symmetric operator.

A forward Euler discretization with step $\Delta t$ is used:
\[
\bm{x}_i(k+1) = \bm{x}_i(k) + \Delta t\, \dot{\bm{x}}_i(k).
\]

\end{document}